% 第13回　プロジェクトの総括（記入例・模範解答）
\session{13}{2}{2027年1月20日（水）}{プロジェクトの総括}{}{}

\goals{
  \item 半年間の作業を、事実、うまくいったこと、うまくいかなかったこと、次に試すことに分けて振り返られる
  \item AIが得意だったこと、頼れなかったこと、人間が判断したことを具体的に挙げられる
  \item 「AIとの協働における人間ならではの役割」を、自分の経験に基づいて説明できる
}

\work[70mm]{個人ワーク1}{AI活用記録を読み返し、印象に残った場面を3つ選ぶ}
\begin{wstabh}{colspec={Q[c,wd=10mm,m] X[2,l,m] X[2,l,m] X[3,l,m]},row{2-Z}={ht=18mm}}
  回 & 場面・AIに頼んだこと & AIの出力 & 自分はどう判断したか、なぜか \\
  \ans{3} & \af{問いに関連する文献を挙げてもらった} & \af{5件のうち2件が実在せず、2件は内容が違った} & \af{AIを「探す場所の見当」に限って使うことにした。実在を確かめる手間は省けないと実感した} \\
  \ans{5} & \af{問いを査読者の立場で批判してもらった} & \af{「過去との比較は半年では難しい」など4つの指摘} & \af{3つは受け入れて問いを修正し、「新しさ」の指摘には反論した。全部を受け入れないことも判断だと気づいた} \\
  \ans{10} & \af{提案の効果を示す根拠を探してもらった} & \af{来店が増えるという具体的な数値を示した} & \af{出典がなかったので使わなかった。数字がなくても、段階的に試す提案にすれば説得できると考えた} \\
\end{wstabh}

\work[100mm]{個人ワーク2}{AIが担ったこと、自分が担ったこと}
\begin{wstabh}{colspec={Q[l,wd=26mm,m] X[l,m] X[l,m]},row{2-Z}={ht=11mm}}
  作業 & AIが担ったこと & 自分（たち）が担ったこと \\
  問いを立てる & \af{6つの視点から問いの候補を出す} & \af{候補を採点し、自分が本当に知りたい問いを選ぶ} \\
  文献を探す・読む & \af{キーワードの提案、要約（読む前の地図）} & \af{データベースで探し、原文で確かめる。要約の誤りを見つける} \\
  問いを確定する & \af{査読者役として弱点を指摘} & \af{どの指摘を受け入れ、どれに反論するかを決める} \\
  ペルソナ・コンセプト & \af{ペルソナ案と30のコンセプト案を出す} & \af{課題に合うペルソナを選び、ペルソナの言葉で案を判定する} \\
  空間・動線 & \af{空間の言語化、画像生成、店内行動の予測} & \af{画像の違和感を言葉にし、つまずきに優先順位を付ける} \\
  提案の検証 & \af{企業側の立場からの批判} & \af{対応・反論・割り切るを決め、一次資料で裏づける} \\
  プレゼン & \af{構成案とコピー案} & \af{見出しを自分たちの言葉で書き、聞き手に合わせて削る} \\
\end{wstabh}

\work[50mm]{振り返りの問い}{自分の経験に基づいて答える}
\begin{wstab}{colspec={X[2,l,m] X[3,l,m]},row{1-Z}={ht=16mm}}
  AIの出力を採用しなかった場面はどこか。なぜそう判断したか & \af{第10回の来店増の数値。出典がなく、企業の担当者に問われたら説明できないから。} \\
  AIに聞かなくても自分で考えたほうが早かった、または良かったのはどんなときか & \af{つまずきの優先順位づけ。ペルソナを一番よく知っているのは、作った自分たちだったから。} \\
  AIがいなければできなかったことは何か & \af{30のコンセプト案を短時間で出すこと。「呼ぶまで話しかけない」という反転の発想は、自分たちだけでは出なかった。} \\
  企業の担当者や他のメンバーとのやりとりで、AIには代われなかったのはどんな役割か & \af{「声をかけないのは現場では勇気がいる」という講評。現場の感覚は、AIの予測からは出てこなかった。} \\
  次に同じようなプロジェクトをするなら、AIの使い方をどう変えるか & \af{自分たちの考えを先に書いてから、AIに批評させる。最初にAIの案を見ると、それに引きずられたから。} \\
\end{wstab}

\begin{promptbox}[自分の記録を分析させる]
次は、私がこの半年に生成AIを使った記録の要約です。AIが役に立った場面と、AIの出力を採用しなかった場面を分類し、そこから読み取れる「私が自分で判断していること」の傾向を指摘してください。AIとしての自己評価ではなく、私の判断の傾向に注目してください。記録：（要約）
\end{promptbox}

\work{共有・討論}{グループでの共通点と違い／AIとの付き合い方はどう変わったか}
\atwoboxes{共通点と違い}{AIとの付き合い方の変化}{32mm}
{共通点：全員が「数値と文献は自分で確かめる」ようになった。違い：メンバーＣは画像生成を多く使い、メンバーＤはほとんど使わず手描きで考えた。}
{最初は「答えを聞く相手」だったが、今は「案を出させて、選ぶのは自分」の相手になった。批判役として使うと、自分では気づかない弱点が見える。}

\work{定義}{AIとの協働における「人間ならではの役割」とは}
\abox[自分の定義（一文で）]{18mm}{AIが広げた選択肢の中から、目の前の相手（顧客や企業）にとって意味のあるものを選び、その選択に責任を持つこと。}
\abox[その根拠となる経験]{36mm}{AIはペルソナもコンセプトも大量に出したが、どれを選ぶかの基準は持っていなかった。私たちは「ペルソナの言葉」と「企業の課題」を基準に選び、選ばなかった理由も説明できた。また、出典のない数値を使わない判断は、企業の担当者に説明する責任があるからこそできた。講評で受けた現場の感覚も、選ぶ基準を磨く材料になった。}

\work[30mm]{提出物}{振り返りレポートの要件チェック}
\begin{achecks}
  \item 「人間ならではの役割」についての自分の定義
  \item その根拠となる経験
\end{achecks}
\afillin{書式・締切（授業内で案内）}{（授業内で案内した書式・締切を記入）}

\begin{point}
  \item 振り返りは、AI活用記録という事実に基づいているかを見る。「AIは便利だった」のような一般論ではなく、何回目のどの場面かが書けていることを求める。
  \item 「人間ならではの役割」には正解がない。評価するのは、定義と根拠となる経験が結びついているか。記入例の定義は一例にすぎない。
  \item 「AIの出力を採用しなかった場面」が書けない学生は、記録が不十分だった可能性がある。記録の取り方を次の機会に活かすよう伝える。
  \item 全体討論では、AIを使わなかった学生の判断も尊重する。使わないことも一つの選択として扱う。
\end{point}
